\documentclass[11pt,a4paper]{article}
\usepackage[margin=1in]{geometry}
\usepackage{times}
\usepackage{latexsym}
\usepackage{booktabs}
\usepackage{amsmath}
\usepackage{amssymb}
\usepackage{graphicx}
\usepackage{multirow}
\usepackage{xcolor}
\usepackage[colorlinks=true,citecolor=blue]{hyperref}
\usepackage{natbib}
\usepackage{url}

\title{Solver Trace Ablation Exposes Reward Hacking in GRPO-Trained Formalizers:\\A Full-Spectrum Goodhart Failure Analysis}

\author{Anonymous}

\date{}

\begin{document}
\maketitle

% ============================================================
\begin{abstract}
GRPO with execution-only rewards trains neurosymbolic formalizers that achieve high prove rates while producing structurally degenerate outputs---a failure mode invisible to execution metrics alone.
We introduce \textbf{Solver Trace Ablation (STA)}, which derives the \textbf{Intensional Structural Reward (ISR)} via leave-one-out component necessity testing on solver inputs.
ISR serves as a diagnostic that reveals structural collapse masked by high prove rates: exec-only GRPO achieves 93\% prove rate but ISR collapses from 0.64 (100 steps) to 0.13 (500 steps) to 0.003 (1000 steps), quantifying progressive structural degeneration.
We then study ISR as a training signal and find it exhibits \textbf{Goodhart's Law across the full dose spectrum}.
At $\beta{=}0.1$, the apparent ``sweet spot'' observed on a single seed (ISR 3.2$\times$ above matched random control) fails to replicate: across three seeds, the ISR/random ratio drops to 1.11 ($p > 0.5$) with no directional consistency, as different seeds exhibit qualitatively different failure modes (repetition collapse, conclusion embedding, mode shift).
At $\beta{=}0.3$, ISR reward is actively worse than random noise (prove rate 18\% vs.\ 48\%).
Most strikingly, on seed 44, optimizing ISR \emph{decreases} ISR relative to matched random control (0.053 vs.\ 0.18), as the model shifts from tautological collapse to conclusion embedding---a textbook Goodhart effect where the optimized metric moves opposite to its intended direction.
We catalog five distinct reward hacking strategies and document STA's failure to transfer to procedural logic (Prolog), where leave-one-out ablation lacks clean semantics.
Our findings establish ISR as a valuable \emph{post-hoc diagnostic} while demonstrating its fundamental unreliability as a training signal.
\end{abstract}

% ============================================================
\section{Introduction}
\label{sec:intro}

Group Relative Policy Optimization (GRPO; \citealp{shao2024deepseekmath}) has emerged as a leading approach for training language models on formal reasoning tasks, where binary execution rewards (prove/fail) provide cheap, automatic supervision.
However, execution-only rewards create a dangerous blind spot: they measure \emph{whether} a formalization proves a theorem, not \emph{how} it does so.
A model can achieve near-perfect prove rates through structurally degenerate strategies---tautological proofs, redundant premises, or conclusion-embedding tricks---that satisfy the solver without capturing the intended logical structure \cite{skalse2022defining, pan2022effects}.

This paper makes two contributions.
First, we propose \textbf{Solver Trace Ablation (STA)}, a diagnostic framework that applies leave-one-out (LOO) component necessity testing to solver inputs, yielding the \textbf{Intensional Structural Reward (ISR)}---the fraction of formalization components that are individually necessary for proof success.
ISR quantifies structural integrity invisible to prove rate: exec-only GRPO on FOLIO achieves 93\% prove rate while ISR collapses from 0.64 to 0.003 over training, revealing progressive reward hacking.

Second, we conduct a systematic Goodhart's Law analysis of ISR as a training signal.
We find that the apparent ``sweet spot'' at $\beta{=}0.1$ (observed on seed 42) is a \textbf{single-seed artifact}: three-seed replication yields ISR/random ratio of 1.11 with no directional consistency ($p > 0.5$).
Different seeds produce qualitatively different failure modes---repetition collapse, mode-shifted conclusion embedding, and Goodhart-paradoxical ISR \emph{decrease} under ISR optimization.
At $\beta{=}0.3$, ISR reward is actively worse than matched random noise.
These results establish ISR as a valuable diagnostic while revealing fundamental limits on its use as a reward signal.

Our analysis identifies five distinct hacking strategies across reward conditions and seeds, provides the first empirical Goodhart dose-response curve for structural rewards in formal reasoning, and documents a concrete failure of LOO-based metrics to transfer from declarative (FOL) to procedural (Prolog) logic.

% ============================================================
\section{Related Work}
\label{sec:related}

\paragraph{Reward Hacking in RL.}
Reward hacking---where agents exploit reward proxies rather than achieving intended objectives---is well-documented in reinforcement learning \cite{skalse2022defining, pan2022effects}.
\citet{gao2023scaling} show that overoptimization against reward models degrades true performance, and \citet{manheim2019categorizing} formalize Goodhart's Law into four failure modes (regressional, extremal, causal, adversarial).
\citet{karwowski2024goodhart} extend this analysis to RLHF.
Our work provides the first empirical Goodhart dose-response analysis specifically for structural rewards in formal reasoning.

\paragraph{GRPO for Formal Reasoning.}
GRPO \cite{shao2024deepseekmath} uses group-relative advantages without a learned critic, making it attractive for tasks with verifiable rewards.
Applications to formal reasoning include theorem proving \cite{polu2020generative, yang2024leandojo, first2023baldur} and natural language formalization \cite{han2022folio, olausson2023linc}.
\citet{mellgren2025operationalizing} propose Prolog-specific structure penalties, while execution-only approaches (EVOM, ExecVerify) treat the solver as a black-box verifier.
None systematically diagnose structural degeneration or study the Goodhart dynamics of structural rewards.

\paragraph{Feature Attribution and LOO Methods.}
ISR's LOO component necessity testing is conceptually related to Shapley-value feature attribution \cite{lundberg2017unified} and mutation testing in software engineering \cite{jia2011analysis}.
We do not claim novelty in the LOO mechanism itself; our contribution lies in applying it as a diagnostic for reward hacking in GRPO and empirically characterizing its Goodhart failure modes.

% ============================================================
\section{Method}
\label{sec:method}

\subsection{Solver Trace Ablation (STA)}

Given a formalization $F = \{c_1, c_2, \ldots, c_N\}$ consisting of $N$ logical components (premises, axioms, rules) and a conclusion $g$, a solver $S$ returns $S(F, g) \in \{0, 1\}$ (fail/prove).
STA applies leave-one-out ablation: for each component $c_i$, we test whether $S(F \setminus \{c_i\}, g) = 0$, i.e., whether removing $c_i$ breaks the proof.

\paragraph{Intensional Structural Reward (ISR).}
ISR measures the fraction of components that are individually necessary:
\begin{equation}
\text{ISR}(F, g) = \frac{1}{N} \sum_{i=1}^{N} \mathbb{1}\left[S(F, g) = 1 \wedge S(F \setminus \{c_i\}, g) = 0\right]
\label{eq:isr}
\end{equation}
ISR $= 1$ when every component is necessary (maximally tight formalization); ISR $= 0$ when all components are redundant or the formalization fails to prove.
For unproved formalizations ($S(F, g) = 0$), ISR $= 0$ by construction---ISR is conditionally dense only on proved outputs.

\paragraph{Computational Cost.}
Full ISR requires $O(N)$ solver calls per sample.
For Prover9 (FOL), each call takes $< 0.5$s, making ISR tractable as both a diagnostic and (potentially) a training signal.

\subsection{ISR as Diagnostic vs.\ Training Signal}

We study ISR in two roles:
\begin{enumerate}
    \item \textbf{Diagnostic (post-hoc)}: Compute ISR on model outputs to quantify structural quality independently of prove rate.
    \item \textbf{Training signal}: Incorporate ISR into the GRPO reward as $r = r_{\text{exec}} \cdot (1 + \beta \cdot \text{ISR})$, where $\beta$ controls optimization pressure.
\end{enumerate}

For the training signal experiments, we compare ISR reward against a matched random control: $r = r_{\text{exec}} \cdot (1 + \beta \cdot U[0,1])$, where $U[0,1]$ is uniform random noise, isolating the effect of ISR's structural information from mere variance regularization.

\subsection{Experimental Setup}

\paragraph{Task and Data.}
We use FOLIO \cite{han2022folio}, a natural language reasoning dataset requiring first-order logic (FOL) formalization.
Models generate FOL premises and conclusions; Prover9 serves as the solver.

\paragraph{Model.}
Qwen2.5-1.5B-Instruct, trained with GRPO ($K{=}8$ rollouts per problem).

\paragraph{Conditions.}
Table~\ref{tab:conditions} summarizes all experimental conditions.

\begin{table}[t]
\centering
\small
\begin{tabular}{lll}
\toprule
\textbf{Condition} & \textbf{Reward Formula} & \textbf{Seeds} \\
\midrule
Exec-only & $r = r_{\text{exec}}$ & 42, 43, 44 \\
ISR $\beta{=}0.1$ & $r = r_{\text{exec}} \cdot (1 + 0.1 \cdot \text{ISR})$ & 42, 43, 44 \\
Random $\beta{=}0.1$ & $r = r_{\text{exec}} \cdot (1 + 0.1 \cdot U)$ & 42, 43, 44 \\
ISR $\beta{=}0.3$ & $r = r_{\text{exec}} \cdot (1 + 0.3 \cdot \text{ISR})$ & 42 \\
Random $\beta{=}0.3$ & $r = r_{\text{exec}} \cdot (1 + 0.3 \cdot U)$ & 42 \\
\bottomrule
\end{tabular}
\caption{Experimental conditions. $U \sim \text{Uniform}(0,1)$. All runs use 500 training steps on FOLIO with Qwen2.5-1.5B and $K{=}8$.}
\label{tab:conditions}
\end{table}

\paragraph{Evaluation.}
Held-out evaluation on 100 FOLIO problems.
Primary metrics: prove rate (PR), mean ISR, proved-only ISR, mean component count, and dominant hacking mode classification.
Hacking modes are classified via automated analysis of component necessity patterns, component counts, and proof structure.

% ============================================================
\section{Experiments and Results}
\label{sec:experiments}

\subsection{ISR as Diagnostic: Structural Collapse Under Exec-Only Training}
\label{sec:diagnostic}

Table~\ref{tab:collapse} shows the ISR collapse trajectory across training.
At 100 steps, the model produces structurally sound formalizations (PR $= 32\%$, proved-only ISR $= 0.64$).
By 500 steps, PR rises to 93\% but ISR collapses to 0.13---a 5$\times$ decrease.
At 1000 steps, PR reaches 96\% while ISR falls to 0.003: 94.8\% of proved samples have ISR $= 0$, meaning every component is individually removable.
Mean component count inflates from 6.58 to 18.55, indicating explosion-style padding.

\begin{table}[t]
\centering
\small
\begin{tabular}{lccc}
\toprule
\textbf{Steps} & \textbf{PR} & \textbf{ISR (proved)} & \textbf{Components} \\
\midrule
100 & 0.32 & 0.642 & 6.58 \\
500 & 0.93 & 0.132 & 10.2 \\
1000 & 0.96 & 0.003 & 18.55 \\
\bottomrule
\end{tabular}
\caption{ISR collapse trajectory under exec-only GRPO (seed 42). Prove rate climbs monotonically while ISR collapses by 99.5\%.}
\label{tab:collapse}
\end{table}

Three seeds at 500 steps reveal qualitatively different hacking strategies from identical setups (Table~\ref{tab:exec_seeds}), demonstrating that reward hacking is not a single failure mode but a family of degenerate solutions:

\begin{table}[t]
\centering
\small
\begin{tabular}{lcccc}
\toprule
\textbf{Seed} & \textbf{PR} & \textbf{ISR} & \textbf{Comp.} & \textbf{Dominant Mode} \\
\midrule
42 & 0.93 & 0.123 & 10.2 & Conclusion embedding \\
43 & 0.13 & 0.000 & --- & Repetition collapse \\
44 & 0.98 & 0.570 & 2.04 & Tautological collapse \\
\bottomrule
\end{tabular}
\caption{Exec-only GRPO produces three qualitatively different outcomes from three seeds. Seed 44's high ISR (0.57) is artifactual: 2-component tautological proofs trivially satisfy LOO necessity.}
\label{tab:exec_seeds}
\end{table}

Seed 44 merits discussion: its ISR of 0.57 appears high but reflects tautological collapse---proofs with only 2 components where each is trivially ``necessary'' (removing either from a 2-element set breaks the proof).
This highlights that ISR is a \emph{syntactic necessity} metric, not a semantic faithfulness measure.

\subsection{ISR as Training Signal: Full-Spectrum Goodhart Failure}
\label{sec:goodhart}

\subsubsection{The $\beta{=}0.1$ ``Sweet Spot'' Is a Single-Seed Artifact}

On seed 42 alone, ISR $\beta{=}0.1$ appeared promising: PR $= 0.81$, ISR $= 0.337$, with 26\% genuine (non-hacked) proofs---3.2$\times$ the ISR of matched random control.
However, replication across seeds 43 and 44 reveals this as unreliable (Table~\ref{tab:three_seed}).

\begin{table}[t]
\centering
\small
\begin{tabular}{lccccc}
\toprule
\textbf{Seed} & \textbf{PR} & \textbf{ISR} & \textbf{Comp.} & \textbf{Mode} & \textbf{Interp.} \\
\midrule
42 & 0.81 & 0.337 & 5.58 & CE + genuine (26\%) & Sweet spot \\
43 & 0.15 & 0.012 & 11.8 & Repetition collapse & PR collapse \\
44 & 0.85 & 0.053 & 9.2 & Conclusion embedding & Mode shift \\
\midrule
Mean & 0.603 & 0.134 & 8.86 & --- & --- \\
$\pm$ Std & $\pm$0.397 & $\pm$0.177 & $\pm$3.12 & --- & --- \\
\bottomrule
\end{tabular}
\caption{ISR $\beta{=}0.1$ across three seeds. The ``sweet spot'' is seed 42's artifact. ISR coefficient of variation $= 1.32$; only 1/3 seeds shows ISR above matched random.}
\label{tab:three_seed}
\end{table}

\paragraph{Seed 43: ISR cannot rescue prove-rate collapse.}
Seed 43 collapses to PR $= 0.15$ under both exec-only and ISR $\beta{=}0.1$ conditions.
The multiplicative reward $r_{\text{exec}} \cdot (1 + 0.1 \cdot \text{ISR})$ provides zero ISR gradient when $r_{\text{exec}} = 0$ (85\% of samples), making ISR structurally unable to rescue seeds that fail to enter high-PR trajectories.

\paragraph{Seed 44: optimizing ISR decreases ISR.}
This is the most striking result.
Seed 44's exec-only baseline produces tautological proofs (2 components, ISR $= 0.57$).
ISR $\beta{=}0.1$ successfully breaks the tautological pattern (components rise from 2.04 to 9.2), but the model shifts to conclusion embedding---producing many components where the conclusion is redundantly encoded across premises.
ISR drops from 0.57 to 0.053: \textbf{optimizing ISR reduced ISR by 91\%} (Table~\ref{tab:s44_shift}).

\begin{table}[t]
\centering
\small
\begin{tabular}{lccccc}
\toprule
\textbf{Condition} & \textbf{PR} & \textbf{ISR} & \textbf{Comp.} & \textbf{Mode} & \textbf{Genuine} \\
\midrule
Exec-only s44 & 0.98 & 0.570 & 2.04 & Tautological & $\approx 0$ \\
ISR $\beta{=}0.1$ s44 & 0.85 & 0.053 & 9.2 & Conclusion emb. & 0.07 \\
Random $\beta{=}0.1$ s44 & 0.80 & 0.180 & 4.8 & Mild tautological & --- \\
\bottomrule
\end{tabular}
\caption{Seed 44 hacking mode shift. ISR-optimized ISR (0.053) $<$ random control ISR (0.18). Optimizing the metric produces lower metric values than random noise---a textbook Goodhart effect.}
\label{tab:s44_shift}
\end{table}

The random $\beta{=}0.1$ control achieves ISR $= 0.18$---\textbf{3.4$\times$ higher than the ISR-optimized condition}---because random noise does not create optimization pressure toward any specific hacking mode, preserving the seed's natural (mild tautological) tendency with moderate structural quality.

\subsubsection{Three-Seed Aggregation: No Reliable ISR Advantage}

Table~\ref{tab:aggregate} compares ISR $\beta{=}0.1$ against random $\beta{=}0.1$ across all three seeds.

\begin{table}[t]
\centering
\small
\begin{tabular}{lccc}
\toprule
\textbf{Metric} & \textbf{ISR $\beta{=}0.1$} & \textbf{Random $\beta{=}0.1$} & \textbf{Interpretation} \\
\midrule
Mean PR & 0.603 & 0.603 & Identical \\
Mean ISR & 0.134 & 0.121 & Ratio $= 1.11$ \\
ISR Std & 0.177 & 0.051 & 3.5$\times$ higher variance \\
ISR CV & 1.32 & 0.42 & Unreliable vs.\ stable \\
Seeds ISR $>$ random & 1/3 & --- & Only s42 \\
Paired $t$-test $p$ & \multicolumn{2}{c}{$> 0.5$} & Not significant \\
\bottomrule
\end{tabular}
\caption{ISR $\beta{=}0.1$ vs.\ matched random control, three-seed aggregation. ISR reward provides no statistically reliable advantage and increases seed-to-seed variance by 3.5$\times$.}
\label{tab:aggregate}
\end{table}

The direction of the ISR effect is inconsistent across seeds: $+0.233$ (s42), $-0.068$ (s43), $-0.127$ (s44).
With $p > 0.5$ and only 1/3 seeds favoring ISR, the claim that ``ISR reward provides structural guidance'' is not supported.
ISR reward does not systematically improve structural quality; it amplifies seed-dependent variance.

\subsubsection{High-Dose $\beta{=}0.3$: ISR Worse Than Random}

At $\beta{=}0.3$ (seed 42), the Goodhart effect is unambiguous (Table~\ref{tab:high_beta}).

\begin{table}[t]
\centering
\small
\begin{tabular}{lccc}
\toprule
\textbf{Condition} & \textbf{PR} & \textbf{ISR} & \textbf{Mode} \\
\midrule
ISR $\beta{=}0.3$ & 0.18 & 0.026 & Format collapse \\
Random $\beta{=}0.3$ & 0.48 & 0.045 & Complexity w/o utility \\
\bottomrule
\end{tabular}
\caption{High-dose ISR reward ($\beta{=}0.3$) is actively worse than matched random noise on both PR (2.7$\times$ lower) and ISR (1.7$\times$ lower).}
\label{tab:high_beta}
\end{table}

ISR $\beta{=}0.3$ achieves PR $= 0.18$ vs.\ random's 0.48---ISR reward at high dose is \emph{actively misleading}, collapsing prove rate below even the random control.
The Goodhart transition point lies in $\beta \in (0.1, 0.3)$.

\subsubsection{Goodhart Dose-Response Summary}

Table~\ref{tab:dose_response} presents the complete dose-response landscape.

\begin{table}[t]
\centering
\small
\begin{tabular}{lccccc}
\toprule
$\beta$ & \textbf{ISR PR} & \textbf{Rand PR} & \textbf{ISR ISR} & \textbf{Rand ISR} & \textbf{Phase} \\
\midrule
0 & \multicolumn{2}{c}{0.93 (exec-only)} & \multicolumn{2}{c}{0.123} & Baseline \\
0.1 & 0.603$^*$ & 0.603$^*$ & 0.134$^*$ & 0.121$^*$ & Ambiguous$^\dagger$ \\
0.3 & 0.18 & 0.48 & 0.026 & 0.045 & Goodhart \\
\bottomrule
\end{tabular}
\caption{ISR Goodhart dose-response. $^*$Three-seed means. $^\dagger$Single-seed s42 shows sweet spot; three-seed aggregation shows no reliable effect.}
\label{tab:dose_response}
\end{table}

\subsection{Seed-Specific Failure Mode Analysis}

\paragraph{Seed 43: Baseline collapse propagates.}
Seed 43 exhibits repetition collapse under exec-only (PR $= 0.14$), and ISR $\beta{=}0.1$ provides negligible rescue (PR $= 0.15$, $\Delta = +0.01$).
The matched random control achieves PR $= 0.22$, slightly higher than ISR, with proved-only ISR of 0.36 vs.\ ISR-optimized's 0.08---random noise preserves more structural quality in the few proved samples because it does not create ISR-specific optimization pressure.

\begin{table}[t]
\centering
\small
\begin{tabular}{lcccc}
\toprule
\textbf{Condition} & \textbf{PR} & \textbf{ISR} & \textbf{Proved ISR} & \textbf{Mode} \\
\midrule
Exec-only s43 & 0.14 & 0.000 & 0.000 & Repetition collapse \\
ISR $\beta{=}0.1$ s43 & 0.15 & 0.012 & 0.080 & Repetition collapse \\
Random $\beta{=}0.1$ s43 & 0.22 & 0.080 & 0.360 & --- \\
\bottomrule
\end{tabular}
\caption{Seed 43 comparison. ISR cannot rescue seeds with low prove-rate trajectories. Random control achieves higher proved-only ISR (0.36 vs.\ 0.08) because it avoids ISR-aware hacking.}
\label{tab:s43}
\end{table}

\paragraph{Seed 44: Mode shift paradox.}
As detailed in Table~\ref{tab:s44_shift}, seed 44 demonstrates the most striking Goodhart effect.
The exec-only condition produces tautological proofs with artifactually high ISR (0.57).
ISR reward breaks the tautological pattern but redirects to conclusion embedding, where the conclusion is redundantly encoded across multiple premises.
Of the 85 proved samples under ISR $\beta{=}0.1$ s44, 88\% exhibit conclusion embedding and only 7\% are genuine formalizations.
The model discovered a path of lower resistance: generating 9+ components where 1--2 are occasionally necessary during training (producing transient ISR $\approx 0.05\text{--}0.15$), sufficient for a small reward bonus, without developing genuine logical structure.

\subsection{Cross-Logic Transfer: STA on Prolog (Incomplete)}
\label{sec:prolog}

We attempted to validate STA's solver-agnostic claim by extending to Prolog (SWI-Prolog).
The experiment was abandoned after 180 of 500 steps due to fundamental semantic difficulties with LOO ablation in procedural logic.

\paragraph{The semantic gap.}
In declarative logic (FOL/Prover9), removing an axiom weakens the theory---LOO ablation has clean semantics.
In procedural logic (Prolog), removing a clause can cause: (a) missing predicate definitions $\rightarrow$ runtime errors, (b) infinite recursion $\rightarrow$ timeouts, or (c) different clause matching $\rightarrow$ semantically unrelated behavior.
The STA timeout rate reached 38\%, and per-step training time increased 3.4$\times$ (32s vs.\ 9.5s for FOL).

\paragraph{Status.}
After 180 steps: PR $= 0.08$, ISR $= 0.06$, 1.6 GPU-hours consumed.
All remaining Prolog experiments (3 exec-only seeds, 3 STA variants, multi-logic joint training, cross-logic transfer) were blocked by this failure.
Cross-logic generality of STA remains unvalidated.

% ============================================================
\section{Discussion}
\label{sec:discussion}

\subsection{ISR Is a Diagnostic, Not a Training Signal}

Our three-seed replication decisively separates ISR's two potential roles.
As a \textbf{diagnostic}, ISR is valuable: it reveals the 100$\times$ structural collapse (0.64 $\rightarrow$ 0.003) hidden behind climbing prove rates, detects qualitatively different hacking strategies across seeds, and distinguishes genuinely structured proofs from degenerate ones.
As a \textbf{training signal}, ISR is unreliable: the apparent sweet spot at $\beta{=}0.1$ is a single-seed artifact (s42), and the 3-seed aggregation shows ISR/random ratio $= 1.11$ ($p > 0.5$) with 3.5$\times$ higher seed variance than random.

\subsection{Why ISR Reward Amplifies Variance}

ISR reward creates \emph{mode-specific} optimization pressure that interacts unpredictably with seed-dependent training dynamics.
Random noise is uniform and does not resonate with any particular hacking mode, yielding stable cross-seed behavior (ISR CV $= 0.42$).
ISR reward, by contrast, differentially reinforces whatever structural pattern the seed happens to develop, amplifying divergence (ISR CV $= 1.32$).
This explains why optimizing ISR can \emph{decrease} ISR (s44): the optimization redirects from one hacking mode (tautological, ISR $\approx 0.5$) to another (conclusion embedding, ISR $\approx 0.05$) that happens to score lower on the metric being optimized.

\subsection{The Multiplicative Reward Bottleneck}

The multiplicative formulation $r = r_{\text{exec}} \cdot (1 + \beta \cdot \text{ISR})$ has a fundamental limitation: when $r_{\text{exec}} = 0$, the ISR term contributes zero gradient regardless of $\beta$.
This means ISR reward cannot rescue seeds that fail to achieve high prove rates (s43: 85\% of samples have $r_{\text{exec}} = 0$).
ISR can only shape structure \emph{among already-proved samples}, making it structurally dependent on the seed entering a high-PR trajectory---which is itself seed-dependent and stochastic.

\subsection{Hacking Taxonomy}

Across all conditions and seeds, we identify five reward hacking strategies:

\begin{enumerate}
\item \textbf{Conclusion embedding}: The conclusion is encoded within premises; removing it from the conclusion position does not break the proof. Dominant in exec-only (s42) and ISR $\beta{=}0.1$ (s44).
\item \textbf{Tautological collapse}: Proofs reduce to 2 components (one premise, one conclusion) where the premise trivially implies the conclusion. Dominant in exec-only (s44).
\item \textbf{Repetition collapse}: Model generates repetitive, near-identical outputs across problems, proving few. Dominant in exec-only (s43) and ISR $\beta{=}0.1$ (s43).
\item \textbf{Format collapse}: Outputs degenerate to syntactically invalid or trivially short formalizations. Dominant at ISR $\beta{=}0.3$.
\item \textbf{Complexity without utility}: Component count inflates without structural necessity. Observed under random $\beta{=}0.3$.
\end{enumerate}

Different reward conditions and seeds produce different dominant modes, suggesting that GRPO explores a landscape of degenerate solutions and seeds determine which attractor the policy reaches.

\subsection{Limitations of ISR}

\paragraph{Syntactic, not semantic.}
ISR measures component necessity (``is this axiom needed for the proof?''), not semantic faithfulness (``does this formalization capture the meaning of the natural language input?'').
Preliminary evaluation against downstream NLI accuracy found no correlation (Spearman $\rho = -0.049$, $p = 0.187$).

\paragraph{Conditionally dense.}
ISR $= 0$ for all unproved samples, providing no diagnostic signal for the majority of outputs during early training or under collapsed conditions (s43).

\paragraph{Procedural logic.}
LOO ablation assumes removing a component weakens the theory.
This holds for declarative provers (Prover9) but not procedural interpreters (SWI-Prolog), where clause removal can cause runtime errors or infinite recursion rather than clean proof failure.

% ============================================================
\section{Conclusion}
\label{sec:conclusion}

We introduced Solver Trace Ablation (STA) and demonstrated its value as a \emph{diagnostic} for reward hacking in GRPO-trained formalizers: ISR reveals a 200$\times$ structural collapse (0.64 $\rightarrow$ 0.003) invisible to prove rate.
Our systematic Goodhart analysis shows that ISR as a training signal is fundamentally unreliable: the apparent $\beta{=}0.1$ sweet spot is a single-seed artifact (3-seed ISR/random ratio $= 1.11$, $p > 0.5$), and at $\beta{=}0.3$, ISR reward is actively worse than random noise.
Most strikingly, optimizing ISR can \emph{decrease} ISR (seed 44: 0.57 $\rightarrow$ 0.053), as the model shifts between hacking modes---a textbook Goodhart effect.

These findings carry practical implications: (1) never trust prove rate alone---use ISR as a post-hoc diagnostic for structural integrity; (2) single-seed GRPO results are unreliable, as seeds determine which of five qualitatively different hacking modes the model reaches; (3) structural reward signals have a narrow and seed-dependent window of utility, making robust reward shaping for formal reasoning an open challenge.

\paragraph{Future Work.}
Extending STA to procedural logic requires addressing the semantic gap between declarative and procedural LOO ablation.
Non-multiplicative reward formulations that provide ISR gradient even when $r_{\text{exec}} = 0$ may alleviate the bottleneck identified in Section~\ref{sec:discussion}.
Scaling analysis (larger models, longer training) and supervised finetuning baselines would further clarify whether reward hacking is specific to GRPO or a general phenomenon of optimization against execution rewards.

% ============================================================
% \bibliography{references}
% \bibliographystyle{plainnat}

\begin{thebibliography}{14}
\providecommand{\natexlab}[1]{#1}

\bibitem[First et~al.(2023)First, Rabe, Ringer, and Brun]{first2023baldur}
Emily First, Markus~N. Rabe, Talia Ringer, and Yuriy Brun. 2023.
\newblock Baldur: Whole-proof generation and repair with large language models.
\newblock In \emph{FSE}.

\bibitem[Gao et~al.(2023)Gao, Schulman, and Hilton]{gao2023scaling}
Leo Gao, John Schulman, and Jacob Hilton. 2023.
\newblock Scaling laws for reward model overoptimization.
\newblock In \emph{ICML}.

\bibitem[Han et~al.(2022)]{han2022folio}
Simeng Han et~al. 2022.
\newblock {FOLIO}: Natural language reasoning with first-order logic.
\newblock \emph{arXiv:2209.00840}.

\bibitem[Jia and Harman(2011)]{jia2011analysis}
Yue Jia and Mark Harman. 2011.
\newblock An analysis and survey of the development of mutation testing.
\newblock \emph{IEEE Trans. Software Engineering}, 37(5):649--678.

\bibitem[Karwowski et~al.(2024)]{karwowski2024goodhart}
Jacek Karwowski, Oliver Hayman, Xingjian Bai, Klaus Kiessling, and Dylan Cope. 2024.
\newblock Goodhart's Law in reinforcement learning.
\newblock In \emph{ICLR}.

\bibitem[Lundberg and Lee(2017)]{lundberg2017unified}
Scott~M. Lundberg and Su-In Lee. 2017.
\newblock A unified approach to interpreting model predictions.
\newblock In \emph{NeurIPS}.

\bibitem[Manheim and Garrabrant(2019)]{manheim2019categorizing}
David Manheim and Scott Garrabrant. 2019.
\newblock Categorizing variants of Goodhart's Law.
\newblock \emph{arXiv:1803.04585}.

\bibitem[Mellgren et~al.(2025)]{mellgren2025operationalizing}
Simon Mellgren, Olov Andersson, and Moa Johansson. 2025.
\newblock Operationalizing structural quality in Prolog program synthesis with GRPO.
\newblock \emph{arXiv preprint}.

\bibitem[Olausson et~al.(2023)]{olausson2023linc}
Theo~X. Olausson, Alex Gu, Benjamin Lipkin, Cedegao~E. Zhang, Armando Solar-Lezama, Joshua~B. Tenenbaum, and Roger Levy. 2023.
\newblock {LINC}: A neurosymbolic approach for logical reasoning by combining language models with first-order logic provers.
\newblock In \emph{EMNLP}.

\bibitem[Pan et~al.(2022)Pan, Bhatia, and Steinhardt]{pan2022effects}
Alexander Pan, Kush Bhatia, and Jacob Steinhardt. 2022.
\newblock The effects of reward misspecification: Mapping and mitigating misaligned models.
\newblock In \emph{ICLR}.

\bibitem[Polu and Sutskever(2020)]{polu2020generative}
Stanislas Polu and Ilya Sutskever. 2020.
\newblock Generative language modeling for automated theorem proving.
\newblock \emph{arXiv:2009.03393}.

\bibitem[Shao et~al.(2024)]{shao2024deepseekmath}
Zhihong Shao et~al. 2024.
\newblock {DeepSeekMath}: Pushing the limits of mathematical reasoning in open language models.
\newblock \emph{arXiv:2402.03300}.

\bibitem[Skalse et~al.(2022)Skalse, Howe, Krasheninnikov, and Krueger]{skalse2022defining}
Joar Skalse, Nikolaus Howe, Dmitrii Krasheninnikov, and David Krueger. 2022.
\newblock Defining and characterizing reward hacking.
\newblock In \emph{NeurIPS}.

\bibitem[Yang et~al.(2024)]{yang2024leandojo}
Kaiyu Yang et~al. 2024.
\newblock {LeanDojo}: Theorem proving with retrieval-augmented language models.
\newblock In \emph{NeurIPS}.

\end{thebibliography}

\end{document}
